%% =============================================================================
%% sm_c4_sandwich.tex — SM-C, Subsection C.4: Sandwich Variance Estimator
%% Body content only (\subsection level); parent structure in sm_c.tex
%% Primary source: dev/log/06_research-notes-survey-design-pseudo-posterior.qmd
%%   lines 1003--1195 (bread/meat matrices, cluster-robust J, SVC absorption)
%% Corrections applied: T2-3 ("exactly" -> "in expectation"), via ref to smb:rem-JeqH
%% =============================================================================

\subsection{Sandwich variance estimator}
\label{smc:sandwich}

The main text (\cref{sec:survey}) defines the bread matrix
$\mathbf{H}_{\mathrm{obs}}$~\eqref{eq:H-obs}, the cluster-robust meat
matrix $\mathbf{J}_{\mathrm{cluster}}$~\eqref{eq:J-cluster}, and the
sandwich variance $\mathbf{V}_{\mathrm{sand}}$~\eqref{eq:V-sand}.  This
subsection provides the detailed derivation: the block-diagonal
structure of~$\mathbf{H}_{\mathrm{obs}}$, the non-block-diagonal
structure of~$\mathbf{J}_{\mathrm{cluster}}$, the cluster-robust
aggregation formula with stratification~\citep{Binder1983}, and the partial absorption of
within-state clustering by the state-varying coefficients.  Throughout,
we use the score functions from~\cref{smb:scores} and the sandwich
convention defined below.

%% ---------------------------------------------------------------------------
%%  C.4.1  Bread matrix H
%% ---------------------------------------------------------------------------

\paragraph{Bread matrix.}
The bread matrix is the negative weighted observed Hessian at the
posterior mean $\hat{\btheta}$,
\begin{equation}\label{smc:eq-H}
  \mathbf{H}_{\mathrm{obs}}
  = -\sum_{i=1}^{N}\tilde{w}_i\,
    \frac{\partial^2\log f_{\HBB}(y_i \mid \btheta)}
         {\partial\btheta\,\partial\btheta\t}
    \bigg|_{\btheta = \hat{\btheta}}.
\end{equation}
By the separability of the hurdle log-likelihood
(\cref{smc:rem-separability-significance}) and the block-diagonal Fisher
information (\cref{smb:prop-block-diag}), the bread inherits a two-block
structure:

\begin{proposition}[Block-diagonal bread matrix]
\label{smc:prop-H-blockdiag}
  The observed Hessian decomposes as
  \begin{equation}\label{smc:eq-H-block}
    \mathbf{H}_{\mathrm{obs}}
    = \begin{pmatrix}
        \mathbf{H}_{\textup{ext}} & \mathbf{0} \\
        \mathbf{0} & \mathbf{H}_{\textup{int}}
      \end{pmatrix},
  \end{equation}
  where\textup{:}
  \begin{enumerate}[label=\textup{(\roman*)},nosep]
    \item The extensive block is the weighted logistic information
    \begin{equation}\label{smc:eq-H-ext}
      \mathbf{H}_{\textup{ext}}
      = \sum_{i=1}^{N}\tilde{w}_i\,q_i(1-q_i)\,
        \mathbf{d}_{1,i}\,\mathbf{d}_{1,i}\t,
    \end{equation}
    with $\mathbf{d}_{1,i}
    = \partial\eta_i^{(1)}/\partial(\balpha\t,\bdelta_{1,s[i]}\t)\t$.
    This matrix is non-stochastic: it depends on $q_i = q_i(\hat{\btheta})$
    but not on the observed data $z_i$.

    \item The intensive block is the weighted negative Hessian of the
    truncated beta-binomial log-likelihood,
    \begin{equation}\label{smc:eq-H-int}
      \mathbf{H}_{\textup{int}}
      = -\sum_{i \in \cl{I}_1}\tilde{w}_i\,
        \frac{\partial^2 \ell_i^{(2+)}}
             {\partial\boldsymbol{\phi}_2\,\partial\boldsymbol{\phi}_2\t}
        \bigg|_{\hat{\btheta}},
    \end{equation}
    where $\boldsymbol{\phi}_2
    = (\bbeta\t,\log\kappa,\bdelta_{2,s[i]}\t)\t$,
    $\cl{I}_1 = \{i : y_i > 0\}$, and
    $\ell_i^{(2+)} = \log f_{\textup{BB}}(y_i \mid n_i,\mu_i,\kappa)
     - \log(1-p_{0,i})$.
    This block depends on the observed counts $y_i$.
  \end{enumerate}
\end{proposition}

\begin{proof}
  By~\cref{smb:prop-block-diag}, the cross-derivative
  $\partial^2\ell_i/\partial\boldsymbol{\phi}_1\,\partial\boldsymbol{\phi}_2\t
  = \mathbf{0}$
  for every observation~$i$.  Multiplying by $\tilde{w}_i > 0$ and
  summing preserves this zero block, giving~\eqref{smc:eq-H-block}.
  The extensive-margin form~\eqref{smc:eq-H-ext} follows from
  $-\partial^2\ell_i^{(1)}/\partial(\eta_i^{(1)})^2 = q_i(1-q_i)$
  (\cref{smb:prop-ext-paramspace}) and the chain rule.
\end{proof}

%% ---------------------------------------------------------------------------
%%  C.4.2  Meat matrix J — individual level and non-block-diagonality
%% ---------------------------------------------------------------------------

\paragraph{Meat matrix and cross-margin coupling.}
The individual-level meat matrix is
\begin{equation}\label{smc:eq-J-ind}
  \mathbf{J}_{\textup{ind}}
  = \sum_{i=1}^{N}\tilde{w}_i^2\,\mathbf{s}_i\,\mathbf{s}_i\t,
\end{equation}
where $\mathbf{s}_i
= \partial\log f_{\HBB}(y_i \mid \btheta)/\partial\btheta
  \big|_{\hat{\btheta}}$
is the observation-level score.

\begin{remark}[Non-block-diagonal meat]
\label{smc:rem-J-nonblockdiag}
  Unlike the bread, the meat matrix is \emph{not} block-diagonal between
  the extensive and intensive margins.  The cross-block is
  \begin{equation}\label{smc:eq-J-cross}
    \mathbf{J}^{\textup{ext,int}}
    = \sum_{i=1}^{N}\tilde{w}_i^2\,
      \mathbf{s}_i^{\textup{ext}}\,
      (\mathbf{s}_i^{\textup{int}})\t.
  \end{equation}
  For $y_i = 0$, the intensive score $\mathbf{s}_i^{\textup{int}}
  = \mathbf{0}$, so the cross-product vanishes.  But for $y_i > 0$,
  both $\mathbf{s}_i^{\textup{ext}} = (1-q_i)\mathbf{d}_{1,i}$ and
  $\mathbf{s}_i^{\textup{int}} \neq \mathbf{0}$, so the cross-block
  receives a nonzero contribution from every positive observation.  In
  consequence, the sandwich variance
  $\mathbf{V}_{\mathrm{sand}}
  = \mathbf{H}_{\mathrm{obs}}^{-1}\mathbf{J}\mathbf{H}_{\mathrm{obs}}^{-1}$
  is not block-diagonal even though $\mathbf{H}_{\mathrm{obs}}$ is.  This
  cross-margin coupling is a general feature of sandwich estimators in
  hurdle models: a single PSU contributes weighted scores from both
  margins simultaneously.
\end{remark}

%% ---------------------------------------------------------------------------
%%  C.4.3  E[J] = H under equal weights
%% ---------------------------------------------------------------------------

\begin{remark}[Reduction under equal weights]
\label{smc:rem-equal-weights}
  Under uniform weights $\tilde{w}_i = 1$ and correct model
  specification, the information identity gives
  $\E[\mathbf{J}_{\textup{ind}}] = \mathbf{H}_{\mathrm{obs}}$
  \citep[see][for the pseudo-posterior analogue]{WilliamsSavitsky2021}, so
  $\mathbf{V}_{\mathrm{sand}} \to \mathbf{H}_{\mathrm{obs}}^{-1}$ and
  the sandwich reduces to the model-based variance.  For the extensive
  margin, $\E[\mathbf{J}_{\textup{ext}}] = \mathbf{H}_{\textup{ext}}$
  holds exactly, because the logistic Hessian is non-random (see
  \cref{smb:rem-ext-info}; under unequal survey weights the identity
  breaks because $\mathbf{H}$ involves $\tilde{w}_i$ while
  $\E[\mathbf{J}]$ involves $\tilde{w}_i^2$).  For the intensive margin, the identity holds only
  in expectation and the finite-sample discrepancy between
  $\mathbf{J}_{\textup{int}}$ and $\mathbf{H}_{\textup{int}}$ is
  nonzero.
\end{remark}

%% ---------------------------------------------------------------------------
%%  C.4.4  Cluster-robust meat (Proposition 4.2)
%% ---------------------------------------------------------------------------

\paragraph{Cluster-robust aggregation.}
The individual-level $\mathbf{J}_{\textup{ind}}$ treats observations as
independent.  In the NSECE stratified cluster design
\citep{NSECEProjectTeam2022}, providers within
the same primary sampling unit (PSU) share neighborhood
characteristics, inducing within-PSU correlation.  The cluster-robust
meat matrix replaces individual scores with PSU-level aggregates,
centered within strata.

\begin{definition}[PSU-level score totals]
\label{smc:def-psu-scores}
  For stratum $h = 1,\ldots,H$ and PSU $c = 1,\ldots,C_h$ within
  stratum~$h$, define
  \begin{equation}\label{smc:eq-s-psu}
    \bar{\mathbf{s}}_{hc}
    = \sum_{i \in \textup{PSU}(h,c)} \tilde{w}_i\,\mathbf{s}_i,
  \end{equation}
  where $\textup{PSU}(h,c) = \{i : \texttt{vstratum}[i]=h,\;
  \texttt{vpsu}[i]=c\}$.  The stratum mean is
  \begin{equation}\label{smc:eq-s-stratum}
    \bar{\mathbf{s}}_h
    = \frac{1}{C_h}\sum_{c=1}^{C_h}\bar{\mathbf{s}}_{hc}.
  \end{equation}
\end{definition}

\begin{proposition}[Cluster-robust meat matrix]
\label{smc:prop-cluster-robust}
  Under the stratified cluster design with $H$~strata and
  $C_h$~PSUs in stratum~$h$, the cluster-robust meat matrix is
  \begin{equation}\label{smc:eq-J-cluster}
    \boxed{%
      \mathbf{J}_{\mathrm{cluster}}
      = \sum_{h=1}^{H}\frac{C_h}{C_h-1}
        \sum_{c=1}^{C_h}
        (\bar{\mathbf{s}}_{hc} - \bar{\mathbf{s}}_h)\,
        (\bar{\mathbf{s}}_{hc} - \bar{\mathbf{s}}_h)\t.
    }
  \end{equation}
  This estimator accounts for\textup{:}
  \textup{(i)}~within-PSU correlation, by aggregating scores to PSU
  totals\textup{;}
  \textup{(ii)}~stratification, by centering PSU scores within strata
  so that only within-stratum variation across PSUs contributes\textup{;}
  \textup{(iii)}~unequal cluster sizes, through the weight-aggregated
  PSU totals $\bar{\mathbf{s}}_{hc}$\textup{;}
  \textup{(iv)}~finite-sample degrees of freedom, via the factor
  $C_h/(C_h-1)$.
\end{proposition}

\begin{proof}
  The design-based variance of the weighted score total under
  stratified sampling decomposes as
  \[
    \Var_{\textup{design}}\!\Bigl(\sum_i\tilde{w}_i\,\mathbf{s}_i\Bigr)
    = \sum_{h=1}^{H}
      \Var_h\!\Bigl(\sum_{c=1}^{C_h}\bar{\mathbf{s}}_{hc}\Bigr),
  \]
  because sampling across strata is independent.  Within stratum~$h$,
  PSUs are sampled without replacement from the population of
  $C_h^{\ast}$~PSUs.  Under the ``ultimate cluster'' assumption that
  $C_h / C_h^{\ast} \approx 0$ (negligible PSU sampling fraction),
  the finite-population correction is dropped and the within-stratum
  variance estimator is
  \[
    \widehat{\Var}_h
    = \frac{C_h}{C_h-1}
      \sum_{c=1}^{C_h}
      (\bar{\mathbf{s}}_{hc}-\bar{\mathbf{s}}_h)\,
      (\bar{\mathbf{s}}_{hc}-\bar{\mathbf{s}}_h)\t,
  \]
  the standard unbiased estimator of the variance of a total from
  $C_h$~units.  Summing over strata gives
  $\mathbf{J}_{\mathrm{cluster}}$.
\end{proof}

\begin{remark}[Degrees of freedom]
\label{smc:rem-dof}
  The cluster-robust estimator has
  $\sum_{h=1}^{H}(C_h - 1) = 415 - 30 = 385$ effective degrees of
  freedom.  With $C_h$ ranging from~5 to~36 (median~$\approx 14$), each
  stratum contributes at least~4 degrees of freedom.  This comfortably
  exceeds the rule-of-thumb threshold of $\sum_h(C_h - 1) > 30$ for
  stable sandwich variance estimation.
\end{remark}

\begin{remark}[Individual versus cluster-level $\mathbf{J}$]
\label{smc:rem-J-comparison}
  When there is no within-PSU correlation conditional on covariates,
  $\mathbf{J}_{\mathrm{cluster}}$ and $\mathbf{J}_{\textup{ind}}$
  converge to the same population quantity.  In practice, within-PSU
  correlation inflates $\mathbf{J}_{\mathrm{cluster}}$ relative to
  $\mathbf{J}_{\textup{ind}}$; the parameter-specific ratio
  $\diag(\mathbf{V}_{\mathrm{sand}})
  / \diag(\mathbf{H}_{\mathrm{obs}}^{-1}
           \mathbf{J}_{\textup{ind}}
           \mathbf{H}_{\mathrm{obs}}^{-1})$
  provides an estimate of the cluster design effect.
  The design effect ratio (DER) classification and diagnostics are
  developed in~\cref{smc:der} below.
\end{remark}

%% ---------------------------------------------------------------------------
%%  C.4.5  SVC partial absorption (Proposition 4.3)
%% ---------------------------------------------------------------------------

\begin{proposition}[Partial absorption of clustering by SVCs]
\label{smc:prop-svc-absorption}
  The state-varying coefficients $\bepsilon_s$ partially absorb
  within-state clustering.  For any two providers $i,j$ in the same
  state~$s$ and the same PSU,
  \begin{equation}\label{smc:eq-svc-absorption}
    \E_{\bepsilon_s}\!\bigl[
      \Cov(s_i,\,s_j \mid \bepsilon_s)
    \bigr]
    \;\leq\;
    \Cov(s_i,\,s_j),
  \end{equation}
  with equality if and only if $\Cov(\E[s_i \mid \bepsilon_s],\,
  \E[s_j \mid \bepsilon_s]) = 0$.
\end{proposition}

\begin{proof}
  By the law of total covariance,
  \[
    \Cov(s_i,\,s_j)
    = \E\!\bigl[\Cov(s_i,\,s_j \mid \bepsilon_s)\bigr]
    + \Cov\!\bigl(
        \E[s_i \mid \bepsilon_s],\;
        \E[s_j \mid \bepsilon_s]
      \bigr).
  \]
  The second term is non-negative, so
  $\Cov(s_i,s_j) \geq \E[\Cov(s_i,s_j \mid \bepsilon_s)]$.
  At the posterior mean $\hat{\btheta}$ (which includes state-specific
  $\hat{\bepsilon}_s$), the plug-in conditional covariance is an estimate
  of $\Cov(s_i,s_j \mid \bepsilon_s)$; hence
  $\mathbf{J}_{\mathrm{cluster}}$ evaluated at $\hat{\btheta}$ will
  be smaller than under a model without SVCs.
\end{proof}

\begin{remark}[Residual clustering]
\label{smc:rem-residual-clustering}
  The inequality~\eqref{smc:eq-svc-absorption} shows that SVCs reduce
  but do not eliminate within-PSU correlation.  The residual correlation
  reflects neighborhood-level effects not captured by the five
  state-varying covariates (poverty, race/ethnicity, urbanicity): local
  labor market conditions, municipal zoning, co-located service
  providers, and other spatially structured unobservables.  The sandwich
  correction with cluster-robust~$\mathbf{J}_{\mathrm{cluster}}$ remains
  necessary to account for this residual design effect.
\end{remark}
